\subsection{Richiedere e completare un acquisto}\label{sec:uc-acquisto}
\begin{xltabular}{\textwidth}{| >{\bfseries}p{.22\textwidth} | >{\raggedright\arraybackslash}X |}
\caption{Caso d'uso: Richiedere e completare un acquisto}\label{tab:uc-acquisto}\\
\hline
\rowcolor{gray!15} Campo & \textbf{Contenuto}\\\hline
\endfirsthead
\multicolumn{2}{c}{{\tablename\ \thetable{} -- Continua dalla pagina precedente}}\\
\hline
\rowcolor{gray!15} Campo & \textbf{Contenuto}\\\hline
\endhead
\hline
\multicolumn{2}{r}{\textit{Continua nella pagina successiva...}}\\
\endfoot
\endlastfoot
Nome & Richiedere e completare un acquisto\\\hline
Livello & Obiettivo Utente\\\hline
Descrizione & Il Customer riserva un veicolo in vendita versando un acconto; dopo l'autorizzazione del Salesman, paga il saldo e conclude l'acquisto. Le azioni possono avvenire in sessioni diverse.\\\hline
Pre-condizioni & Customer e Salesman hanno effettuato l'accesso alle rispettive aree; il veicolo esiste ed è in stato \code{SALE}; il Customer dispone dell'importo necessario per l'acconto.\\\hline
Post-condizioni & Dopo l'acconto il veicolo è \code{BLOCKED} e la richiesta è registrata. Dopo il saldo il veicolo è \code{SOLD}, la richiesta è eliminata, la vendita è registrata e gli eventuali test drive futuri sono cancellati. Ciascuna fase ha una propria transazione.\\\hline
Attori & Customer (primario), Salesman\\\hline
Schermate & Passi 1--2: catalogo (figura~\ref{fig:ui-catalogo}); passo 3: richieste Salesman (figura~\ref{fig:ui-salesman-richieste}); passi 4--6: acquisti del Customer (figura~\ref{fig:ui-acquisti}) e conto (figura~\ref{fig:ui-conto}).\\\hline
Test & \code{CheckoutFunctionalTest.completePurchaseWorkflow} verifica il percorso principale su H2; \code{cancelPurchaseWorkflowWithPenalty} verifica l'annullamento. \code{SalesServiceTest} verifica avvio e penale con mock. Autorizzazione assente, scadenza e variazione del prezzo non sono coperti da una prova funzionale dedicata.\\\hline
Flusso Principale &
\begin{enumerate}[leftmargin=*,nosep]
  \item Il Customer seleziona un veicolo in vendita; l'interfaccia calcola l'acconto, pari al 20\% del prezzo visualizzato. Il saldo finale usa invece \code{Vehicle.price}, che può differire dal prezzo visualizzato in presenza di sconti.
  \item Il Customer conferma: il sistema addebita l'acconto, registra la richiesta e blocca il veicolo.
  \item Il Salesman consulta la richiesta e la autorizza; il sistema registra il venditore e fissa la scadenza a un mese.
  \item Il Customer seleziona la richiesta autorizzata e conferma il pagamento del saldo.
  \item Il sistema calcola la differenza tra il prezzo corrente del veicolo e l'acconto già versato, quindi addebita l'eventuale residuo.
  \item Il sistema registra la vendita, elimina la richiesta, pone il veicolo in stato \code{SOLD} e cancella i test drive futuri relativi al veicolo.
\end{enumerate}\\\hline
Flussi Alternativi &
\textbf{2. Veicolo non disponibile o saldo insufficiente:} la richiesta non viene creata e le scritture della fase sono annullate.\newline\newline
\textbf{4. Autorizzazione assente:} l'interfaccia impedisce di avviare il saldo. Il servizio di completamento non verifica autonomamente autorizzazione e scadenza: una chiamata diretta potrebbe aggirare quel controllo.\newline\newline
\textbf{4. Annullamento o scadenza:} la richiesta viene eliminata e il veicolo torna in vendita; il sistema trattiene il 20\% dell'acconto e rimborsa l'80\%. Le richieste scadute sono elaborate all'avvio e durante alcune consultazioni, senza un timer dedicato.\newline\newline
\textbf{5. Prezzo variato:} la richiesta non conserva un prezzo finale bloccato; il residuo è calcolato sul prezzo del veicolo letto al momento del saldo.\\\hline
\end{xltabular}
